\documentclass[conference]{IEEEtran}
\IEEEoverridecommandlockouts
\usepackage{cite}
\usepackage{amsmath,amssymb,amsfonts}
\usepackage{algorithmic}
\usepackage{graphicx}
\usepackage{textcomp}
\usepackage{xcolor}
\def\BibTeX{{\rm B\kern-.05em{\sc i\kern-.025em b}\kern-.08em
    T\kern-.1667em\lower.7ex\hbox{E}\kern-.125emX}}
\begin{document}

\title{NutriGuard-AI: An Age-Aware, Evidence-Grounded Food Label Assistant for Indian Packaged Foods}

\author{\IEEEauthorblockN{1\textsuperscript{st} Nandhakumar}
\IEEEauthorblockA{\textit{Department of Computer Science} \\
\textit{Independent Researcher}\\
India \\
nandhakumar@example.com}
}

\maketitle

\begin{abstract}
Nutritional scoring systems for packaged foods often lack transparency and fail to account for age-specific metabolic thresholds or the degree of food processing. We present NutriGuard-AI, a multimodal, deterministic scoring engine tailored for Indian retail packaged foods. The system extracts label data via vision-language models and applies a deterministic, multi-factor algorithmic framework incorporating Nutrition, Ingredients, Processing, Additives, and NOVA structural penalties. The revised evaluation characterizes behavioral and implementation properties rather than establishing clinical or population-level nutritional validity. Using an evaluation collection of 100 product records corresponding to 99 unique barcodes, we demonstrate robust reference-system alignment with the established NOVA classification (Spearman $\rho = 0.697$). Internal parameter robustness analysis reveals that the primary grade conclusions are relatively stable under $\pm 50\%$ coefficient perturbations. We also provide a structured software-level schema robustness audit and an evidence coverage audit to quantify citation-coverage gaps. Our architecture enforces a strict boundary between deterministic numerical scoring and generative explanation, offering a transparent, age-aware tool to assist consumers in making healthier dietary choices.
\end{abstract}

\begin{IEEEkeywords}
Nutritional Scoring, RAG, Multimodal LLM, Food Processing, Deterministic Algorithm
\end{IEEEkeywords}

\section{Introduction}
Existing nutritional scoring systems provide valuable epidemiological insights but often operate as black boxes for the end consumer. They frequently omit processing penalties and lack age-group specificity. NutriGuard-AI bridges this gap by combining an explicit, deterministic multi-factor algorithm with a Retrieval-Augmented Generation (RAG) explanation layer grounded strictly in WHO and ICMR-NIN dietary guidelines. 

\section{Related Work}
Recent systems emphasize recommendation, intervention, retrieval, or multimodal interpretation, whereas NutriGuard-AI focuses on exposing component-level contributions within a deterministic packaged-food assessment.
\begin{table}[h]
\caption{Comparison of Nutritional AI Systems}
\begin{center}
\begin{tabular}{|p{2cm}|p{2cm}|p{2cm}|p{1.5cm}|}
\hline
\textbf{Work} & \textbf{Architecture} & \textbf{Validation} & \textbf{Domain} \\
\hline
Dindukurthi et al. \cite{ref1} & GraphRAG & Ranking eval & Personalized \\
Wu et al. \cite{ref2} & Hierarchical multi-agent & Clinical eval & Clinical \\
Prasad et al. \cite{ref3} & Multimodal + RAG & NLP/OCR metrics & Nutrition labels \\
Deng \& Tu \cite{ref4} & Multimodal LLM & Recommendation & Diet rec. \\
Wang et al. \cite{ref5} & LLM-RAG + HEI & HEI eval & General \\
Shikara et al. \cite{ref6} & Multi-agent + RAG & Framework eval & Personalized \\
\textbf{NutriGuard-AI (Ours)} & \textbf{Deterministic + Evidence RAG} & \textbf{Ablation, sensitivity, alignment} & \textbf{Indian retail} \\
\hline
\end{tabular}
\end{center}
\end{table}

\section{Methodology}
\subsection{Algorithmic Framework}
The core scoring engine applies the following formula:
\begin{equation}
S_{\mathrm{pre}} = M_{\mathrm{NOVA}}(0.20I + 0.15P + 0.30A) + 0.35N
\end{equation}
The coefficients ($0.20, 0.15, 0.30, 0.35$) are designed parameters reflecting relative importance based on literature, rather than clinically learned optimal weights. The NOVA multiplier ($M_{\mathrm{NOVA}}$) strictly penalizes ultra-processed foods.

\subsection{Evaluation Protocol}
To evaluate the system, we utilized an evaluation collection of 100 Indian retail product records (representing 99 unique barcodes) across 14 categories.
\begin{itemize}
    \item \textbf{E1: Deterministic Software Verification.} Schema robustness and formatting audits.
    \item \textbf{E2: Component and Demographic Sensitivity.} Component influence tracking across age profiles.
    \item \textbf{E3: Parameter Robustness.} Perturbing coefficients by $\pm 50\%$. When one coefficient was perturbed, the resulting weight difference was proportionally redistributed across the remaining coefficients to preserve the unit-sum constraint.
    \item \textbf{E4: Reference-System Alignment.} Calculating continuous score alignment with the NOVA processing framework using a tie-aware Spearman correlation. A major divergence was operationally defined as a product in NOVA 4 receiving Grade A+, A, B+, or B; a product in NOVA 1 receiving Grade D+, D, or E; or a product in NOVA 3 receiving Grade A+ or A.
    \item \textbf{E5: Evidence Audit.} Auditing evidence coverage gaps in the RAG layer prior to generation.
\end{itemize}

\section{Results}
\subsection{Reference-System Alignment (NOVA)}
The NutriGuard-AI continuous score showed a positive monotonic association with inverted NOVA classification ($\rho=0.697$), indicating that the implemented scoring function generally assigns lower scores to products with higher NOVA processing levels. Because NOVA classification is itself incorporated into the NutriGuard-AI scoring function, this analysis represents reference alignment rather than independent external validation. The dataset is heavily imbalanced toward NOVA 4 (80\%). The grade correlation of 0.721 is actually informative: the grade discretization preserves association despite the heavy skew. There were exactly 0 major divergence cases.

\subsection{Parameter Robustness}
The results indicate that the product ranking is relatively stable under the tested coefficient perturbations, although grade assignments are more sensitive near grade boundaries. Perturbing the Nutrition weight by $\pm 50\%$ yielded a max absolute score change of 7, with rank stability remaining high ($\rho = 0.968-0.979$).

\subsection{Evidence Coverage and Explanation Guardrail Audit}
The evidence audit identified incomplete coverage for four nutrient types (calories, saturated fat, cholesterol, caffeine), demonstrating that the evidence-grounding layer remains an active limitation of the prototype. The implementation includes a deterministic fallback when required evidence is unavailable.

\section{Discussion and Limitations}
NutriGuard-AI demonstrates reproducible deterministic behavior, measurable component sensitivity, coefficient-perturbation stability, age-dependent variation, and alignment with the NOVA processing framework on a fixed 100-record collection. These experiments do not establish clinical validity or prove that the numerical score is an optimal measure of nutritional quality.

\begin{table}[h]
\caption{Validation Hierarchy}
\begin{center}
\begin{tabular}{|l|l|}
\hline
\textbf{Evidence level} & \textbf{NutriGuard-AI status} \\
\hline
Software correctness & Evaluated \\
Determinism & Evaluated \\
Component sensitivity & Evaluated \\
Parameter robustness & Evaluated \\
NOVA reference alignment & Evaluated \\
Schema robustness & Evaluated \\
Evidence coverage & Audited \\
OCR accuracy & Not evaluated \\
Human expert agreement & Not evaluated \\
External nutritional validity & Not established \\
Clinical validity & Not established \\
Health outcomes & Not evaluated \\
\hline
\end{tabular}
\end{center}
\end{table}

\section{Conclusion}
NutriGuard-AI provides a stable, transparent alternative for nutritional scoring. The revised evaluation combines deterministic software verification with reference-system alignment, parameter robustness analysis, and structured evidence coverage audits.

\begin{thebibliography}{00}
\bibitem{ref1} Dindukurthi et al., ``GraphRAG for Personalized Nutrition,'' 2024.
\bibitem{ref2} Wu et al., ``Hierarchical multi-agent systems for clinical nutrition,'' 2025.
\bibitem{ref3} Prasad et al., ``Multimodal LLM and RAG for nutrition labels,'' 2025.
\bibitem{ref4} Deng \& Tu, ``Multimodal LLM in dietary recommendation,'' 2024.
\bibitem{ref5} Wang et al., ``LLM-RAG with HEI for general nutrition,'' 2025.
\bibitem{ref6} Shikara et al., ``Multi-agent RAG framework for personalized nutrition,'' 2025.
\end{thebibliography}

\end{document}
